\section{Explicit factorization of the normal forms}\label{sec:explicit}

We now factor the normal forms under the prefix conditions of
\Cref{sec:combinatorics}.  Both constructions build the factorization as a
sum of rank-one terms indexed by \emph{slots} $t \in \{1, \ldots, n\}$: for
$A = LU$,
\begin{equation}\label{eq:rank-one-lu}
  LU = \sum_{t=1}^{n} \ell_t \, u_t^T,
\end{equation}
where $\ell_t$ is column $t$ of $L$ (supported on rows $\ge t$) and $u_t^T$
is row $t$ of $U$ (supported on columns $\ge t$); for $A = LDL^T$,
\begin{equation}\label{eq:rank-one-ldlt}
  LDL^T = \sum_{t=1}^{n} d_t \, \ell_t \ell_t^T .
\end{equation}
The difference between the two problems is now plain: in
\cref{eq:rank-one-lu} each term is an arbitrary rank-one matrix with the
required support, whereas in \cref{eq:rank-one-ldlt} each term is
\emph{symmetric} rank-one.  A single off-diagonal exchange block
$\alpha(e_i e_j^T + e_j e_i^T)$ has rank $2$ and zero diagonal, and cannot
equal a sum of terms each supported in rows and columns $\ge i$ unless
partial sums are allowed to cancel (over any field: the $(i,i)$ entry
forces $d_t \ell_t[i]^2 = 0$ for every term while the $(i,j)$ entry needs
some $d_t \ell_t[i] \neq 0$); this forces the telescoping chains of
\Cref{thm:ldlt-matching}.

\subsection{LU of a matching matrix: greedy slot assignment}

\begin{theorem}\label{thm:lu-matching}
  Let $\Pi$ be a matching matrix satisfying the prefix condition
  $\nu_k \le k$ for all $k$ (\Cref{lem:comb-lu}(c)).  Then there is an
  injection
  \begin{equation}\label{eq:slot-injection}
    \sigma : \M(\Pi) \hookrightarrow \{1, \ldots, n\},
    \qquad
    \sigma(i,j) \le \min(i,j) \quad \text{for every } (i,j) \in \M(\Pi),
  \end{equation}
  computable by a greedy pass, and setting, for each $(i,j) \in \M$ with
  $t = \sigma(i,j)$,
  \[
    \ell_t := e_i,
    \qquad
    u_t := w_{ij} \, e_j,
  \]
  and $\ell_t = u_t = 0$ for slots $t \notin \range(\sigma)$, the matrices
  $L = (\ell_1, \ldots, \ell_n)$ and $U$ with rows $u_1^T, \ldots, u_n^T$
  are lower and upper triangular, respectively, and $L U = \Pi$.
\end{theorem}

\begin{proof}
  \emph{The assignment suffices.}
  Given $\sigma$ as in \cref{eq:slot-injection}, triangularity is immediate:
  $\ell_t = e_i$ has its support at row $i \ge \min(i,j) \ge t$, and $u_t$
  at column $j \ge t$.  Since $\sigma$ is injective, each matched pair
  contributes its own term in \cref{eq:rank-one-lu} and empty slots
  contribute zero, so
  $L U = \sum_{(i,j) \in \M} e_i \, (w_{ij} e_j)^T = \Pi$.  Note that both
  diagonals may carry zeros: slot $t$ has $L[t,t] = 1$ only if
  $i = t$, and $U[t,t] \neq 0$ only if $j = t$.

  \emph{The assignment exists.}
  Process the pairs in order of nondecreasing $m(i,j) := \min(i,j)$, ties
  broken arbitrarily, maintaining the set of free slots (initially
  $\{1, \ldots, n\}$); assign to each pair the largest free slot
  $t \le m(i,j)$ (the argument below is valid for \emph{any} choice of
  free slot $\le m(i,j)$; we fix the largest for definiteness).  Suppose
  the greedy fails at a pair with bound
  $m := m(i,j)$: all slots $1, \ldots, m$ are occupied.  Every occupied
  slot $s \le m$ was assigned to a previously processed pair, whose bound
  was $\le m$ (bounds are nondecreasing along the pass).  Together with the
  current pair, this exhibits $m + 1$ pairs with $\min \le m$, i.e.\
  $\nu_m \ge m + 1$, contradicting the prefix condition.  Hence the greedy
  assigns every pair and produces $\sigma$.
\end{proof}

\begin{remark}[which pairs need rerouting]\label{rem:rerouting}
  Order the pairs by their bound $m(i,j)$.  A diagonal pair $(t,t)$ can
  take its own slot $t$; a pair $(i,j)$ with $i < j$ (an "above" pair)
  would naturally take slot $i$, and a "below" pair $(i,j)$ with $i > j$
  slot $j$.  Conflicts occur exactly at \emph{crossing indices} $t$ that
  are simultaneously the row of an above pair $(t, j)$ and the column of a
  below pair $(i, t)$: one of the two must be rerouted to an earlier free
  slot, whose availability is exactly what the prefix condition
  guarantees --- in the leanest case the free slot comes from an index that
  is both a zero row and a zero column.  The exchange matrix $A_2$ of
  \cref{eq:swap-vs-exchange} is the minimal example: index $2$ is crossing,
  and the pair $(2,3)$ (or $(3,2)$) is rerouted through slot $1$, whose row
  and column are zero.  This is the precise sense in which zero diagonal
  entries are a \emph{resource} for unpivoted factorization.
\end{remark}

\subsection{$LDL^T$ of a symmetric matching matrix: threads}

Throughout this subsection $P$ is a symmetric matching matrix satisfying
$\chi_k \le \zeta_k$ for all $k$ (\Cref{lem:comb-sym}).

\begin{lemma}[thread partition]\label{lem:threads}
  The pairs of $P$ can be partitioned into \emph{threads}
  \begin{equation}\label{eq:thread}
    r < i_1 < j_1 < i_2 < j_2 < \cdots < i_m < j_m,
  \end{equation}
  where $r$ is an isolated zero index, $(i_q, j_q)$ are pairs of $P$, and
  distinct threads use disjoint sets of indices.
\end{lemma}

\begin{proof}
  Scan the indices from left to right, maintaining a set of \emph{free
  threads}.  When an isolated zero index $r$ is encountered, create a new
  free thread rooted at $r$.  When a pair $(i,j)$ starts at index $i$,
  attach it to any free thread, which becomes busy until index $j$; when
  the scan reaches $j$, the thread becomes free again.  (Each index belongs
  to at most one block, so no two events collide at the same index.)  It
  remains to show a free thread is always available when a pair starts at
  $i$.  The threads created before index $i$ number $\zeta_{i-1} = \zeta_i$
  ($i$ is a pair index, not an isolated zero).  The threads busy at that
  moment are those attached to pairs $(i', j')$ with $i' < i < j'$;
  together with the new pair, these are exactly the pairs counted by
  $\chi_i = \#\{(i', j') : i' \le i < j'\}$, so at most $\chi_i - 1$
  threads are busy when the new pair arrives.  Since $\chi_i \le \zeta_i$,
  at least $\zeta_i - (\chi_i - 1) \ge 1$ of the created threads are
  free.  By construction
  each thread's blocks are attached after its root and after the previous
  block's end, giving the interleaving \cref{eq:thread}, and every index is
  used by at most one thread.
\end{proof}

\begin{lemma}[one thread telescopes]\label{lem:telescope}
  Fix a thread \cref{eq:thread} and write $\alpha_q \neq 0$ and $\beta_q$
  for the weights of the pair $(i_q, j_q)$.  Define, backward for
  $q = m, m-1, \ldots, 1$:
  \[
    \delta_{m+1} := 1,
    \quad
    s_{m+1} := 0 \in \F^n,
    \qquad
    v_q :=
    \begin{cases}
      0 & \beta_q \neq 0\\
      1 & \beta_q = 0
    \end{cases},
    \quad
    \delta_q := \delta_{q+1} v_q^2 - \beta_q,
    \quad
    s_q := -\frac{\alpha_q}{\delta_q} e_{i_q}
    + \frac{\delta_{q+1} v_q}{\delta_q} s_{q+1} .
  \]
  Each $\delta_q$ is nonzero ($\delta_q = -\beta_q \neq 0$ in the first
  case, $\delta_q = \delta_{q+1}$ in the second).  Define the columns and
  diagonal entries
  \[
    \ell_r := s_1, \quad d_r := \delta_1;
    \qquad
    \ell_{i_q} := s_q + e_{j_q}, \quad d_{i_q} := -\delta_q;
    \qquad
    \ell_{j_q} := s_{q+1} + v_q e_{j_q}, \quad d_{j_q} := \delta_{q+1} .
  \]
  Then each $\ell_t$ is supported on rows $\ge t$, and
  \begin{equation}\label{eq:thread-sum}
    d_r \ell_r \ell_r^T
    + \sum_{q=1}^{m}
    \bigl( d_{i_q} \ell_{i_q} \ell_{i_q}^T + d_{j_q} \ell_{j_q}
    \ell_{j_q}^T \bigr)
    =
    \sum_{q=1}^{m}
    \Bigl( \alpha_q \bigl( e_{i_q} e_{j_q}^T + e_{j_q} e_{i_q}^T \bigr)
    + \beta_q \, e_{j_q} e_{j_q}^T \Bigr),
  \end{equation}
  i.e.\ the thread's columns reproduce exactly the thread's blocks of $P$.
\end{lemma}

\begin{proof}
  \emph{Support.}  By downward induction, $s_q$ is supported on
  $\{i_q, i_{q+1}, \ldots, i_m\}$.  Hence $\ell_r = s_1$ is supported on
  indices $\ge i_1 > r$; $\ell_{i_q}$ on
  $\{i_q, \ldots, i_m\} \cup \{j_q\}$, all $\ge i_q$; and $\ell_{j_q}$ on
  $\{i_{q+1}, \ldots, i_m\} \cup \{j_q\}$, all $\ge j_q$ since
  $i_{q+1} > j_q$.

  \emph{Telescoping.}  The recursion gives the two identities
  \begin{equation}\label{eq:thread-identities}
    -\delta_q s_q + \delta_{q+1} v_q \, s_{q+1} = \alpha_q e_{i_q},
    \qquad
    \delta_{q+1} v_q^2 - \delta_q = \beta_q .
  \end{equation}
  Expanding the symmetric rank-one terms of one pair,
  \begin{align*}
    d_{i_q} \ell_{i_q} \ell_{i_q}^T + d_{j_q} \ell_{j_q} \ell_{j_q}^T
    &= -\delta_q (s_q + e_{j_q})(s_q + e_{j_q})^T
    + \delta_{q+1} (s_{q+1} + v_q e_{j_q})(s_{q+1} + v_q e_{j_q})^T \\
    &= \bigl( -\delta_q s_q s_q^T + \delta_{q+1} s_{q+1} s_{q+1}^T \bigr)
    + \bigl( -\delta_q s_q + \delta_{q+1} v_q s_{q+1} \bigr) e_{j_q}^T \\
    &\qquad
    + e_{j_q} \bigl( -\delta_q s_q + \delta_{q+1} v_q s_{q+1} \bigr)^T
    + \bigl( \delta_{q+1} v_q^2 - \delta_q \bigr) e_{j_q} e_{j_q}^T \\
    &= \bigl( -\delta_q s_q s_q^T + \delta_{q+1} s_{q+1} s_{q+1}^T \bigr)
    + \alpha_q \bigl( e_{i_q} e_{j_q}^T + e_{j_q} e_{i_q}^T \bigr)
    + \beta_q \, e_{j_q} e_{j_q}^T,
  \end{align*}
  using \cref{eq:thread-identities}.  Summing over $q = 1, \ldots, m$, the
  first brackets telescope to
  $-\delta_1 s_1 s_1^T + \delta_{m+1} s_{m+1} s_{m+1}^T
  = -\delta_1 s_1 s_1^T$, which is cancelled by the root term
  $d_r \ell_r \ell_r^T = \delta_1 s_1 s_1^T$.  This proves
  \cref{eq:thread-sum}.
\end{proof}

\begin{theorem}\label{thm:ldlt-matching}
  Let $P$ be a symmetric matching matrix with $\chi_k \le \zeta_k$ for all
  $k$.  Then there are a lower triangular $S$ and a diagonal $E$ with
  \[
    P = S E S^T .
  \]
\end{theorem}

\begin{proof}
  Partition the pairs into threads (\Cref{lem:threads}) and apply
  \Cref{lem:telescope} to each thread, obtaining columns
  $\ell_t$ and weights $d_t$ for every index $t$ in a thread.  For each
  singleton $t$ set $\ell_t := e_t$ and $d_t := \gamma_t$; for isolated
  zero indices not rooting a thread set $\ell_t := 0$, $d_t := 0$.  Let
  $S := (\ell_1, \ldots, \ell_n)$ and $E := \operatorname{diag}(d_1,
  \ldots, d_n)$.  Every $\ell_t$ is supported on rows $\ge t$, so $S$ is
  lower triangular.  Distinct threads and singletons use disjoint index
  sets, so summing \cref{eq:thread-sum} over threads and adding the
  singleton terms $\gamma_t e_t e_t^T$ reproduces every block of $P$ and
  nothing else:
  $S E S^T = \sum_t d_t \ell_t \ell_t^T = P$.
\end{proof}

\begin{remark}[the real signed case]\label{rem:signed-threads}
  Over $\mathbb{R}$, after the normalization of \Cref{rem:signed}
  ($\alpha_q = 1$, $\beta_q = 0$), the recursion of \Cref{lem:telescope}
  simplifies to $v_q = 1$, $\delta_q = 1$, and
  $s_q = s_{q+1} - e_{i_q} = -(e_{i_q} + e_{i_{q+1}} + \cdots + e_{i_m})$,
  so the thread columns are signed partial sums and
  \[
    d_r = 1, \qquad d_{i_q} = -1, \qquad d_{j_q} = 1 :
  \]
  the diagonal alternates along each thread, and $D$ can be taken with
  entries in $\{0, \pm 1\}$.  The factorization
  \cref{eq:exchange-factorizations} of $A_2$ is exactly the one-thread case
  $r = 1 < i_1 = 2 < j_1 = 3$, obtained by seeding the recursion with
  $\delta_{m+1} = -1$ instead of $1$ (valid whenever all $\beta_q = 0$),
  which flips the sign of every $\delta_q$ and $s_q$.
\end{remark}
